% AY2024 制御工学 第1問〔I〕（転記）
% 出典: 03_制御工学/original/03_制御工学/2024/2024se.pdf
% 要約: 壁-並列(バネ k2・ダンパ d)-質量 m-バネ k1-入力端 A(変位 x1)。
%       運動方程式・伝達関数、ステップ応答・最大値、(5)で A 固定し質量に外力 f(t)。

\section*{第1問〔I〕}

図1に示す系に関する次の問い (1)〜(5) に答えよ. なお, 質量を $m$, 粘性係数を $d$,
バネ定数を $k_1$, $k_2$ とし, $x_1(t)$ は系への入力となる変位, $x_2(t)$ は質量の変位を表す.
また, 初期状態において系は静止しており, 粘性に関しては, ダッシュポットの動作速度に
比例した粘性抵抗が生じるものとする.

\begin{center}
\begin{tikzpicture}[>=Latex]
  % 壁
  \draw[thick] (0,0.3) -- (0,2.0);
  \foreach \y in {0.4,0.55,...,1.95} \draw[thick] (0,\y) -- (-0.22,\y-0.22);
  % バネ k2（上）
  \draw[thick,decorate,decoration={zigzag,segment length=7,amplitude=4}] (0,1.55) -- (1.6,1.55);
  \node at (0.8,1.95) {$k_2$};
  % ダンパ d（下）
  \draw[thick] (0,0.85) -- (0.45,0.85);
  \draw[thick] (0.45,0.63) rectangle (0.9,1.07);
  \draw[thick] (0.7,0.85) -- (1.6,0.85);
  \node at (0.8,0.35) {$d$};
  % 質量 m
  \draw[thick,fill=gray!12] (1.6,0.45) rectangle (2.35,1.95);
  \node at (1.975,1.2) {$m$};
  \draw[->,thick] (1.975,0.2) -- (2.6,0.2);
  \node at (2.7,-0.02) {$x_2(t)$};
  % バネ k1
  \draw[thick,decorate,decoration={zigzag,segment length=7,amplitude=4}] (2.35,1.2) -- (3.8,1.2);
  \node at (3.075,1.6) {$k_1$};
  % 入力端 A
  \draw[thick] (3.8,1.2) -- (4.2,1.2);
  \draw[thick,fill=white] (4.2,1.2) circle (0.06);
  \node at (4.5,1.2) {A};
  \draw[->,thick] (4.2,0.95) -- (4.85,0.95);
  \node at (4.7,0.72) {$x_1(t)$};
  \node at (2.2,-0.45) {\small 図1};
\end{tikzpicture}
\end{center}

\begin{enumerate}[label=(\arabic*)]
  \item 系の運動方程式を求めたうえで, 入力を $x_1(t)$, 出力を $x_2(t)$ として伝達関数を求めよ.

  \item 各パラメータを $m=1$, $d=2$, $k_1=3$, $k_2=2$ とし, 入力として $x_1(t)$ にステップ
        入力を与えたときの応答 $x_2(t)$ を求めよ. また, 定常値が存在する場合には, 定常値に
        ついても求めよ.

  \item 問い (2) で求めた応答 $x_2(t)$ の最大値とそのときの $t$ の値を求めよ.

  \item 各パラメータを $m=1$, $d=0$, $k_1=3$, $k_2=1$, 初期値をすべて $0$ とし, 入力として
        $x_1(t)$ にステップ入力を与えたときの応答 $x_2(t)$ を求めよ.

  \item 次に, すべてのバネが自然長の状態において点 A を固定端に取付け, 各パラメータを
        $m=1$, $d=2$, $k_1=1$, $k_2=1$ とし, 質量にステップ入力となる力 $f(t)$ を加える.
        この条件下において, 初期値をすべて $0$ としたときの応答 $x_2(t)$ を求めよ.
\end{enumerate}
